\chapter{Discussion}
\label{ch:discussion}
\begin{spacing}{1.4}

\section{Prediction validity before a fairness interpretation}

The findings answer the research questions in a sequence: establish a useful predictor,
check the proxy mechanism, inspect transfer to audited groups, and distinguish the
independent corpus's evidence from a replication of mitigation. The original neural
experiment alone cannot support a broad conclusion about proxy regularisation because
its ordinary model is weaker than the simple references. Its pooled three-seed SRCC of
0.170 and original-scale MAE of 0.951 compare unfavourably with the fixed all-modality
ridge MAE of 0.535. This is an important baseline diagnosis, not evidence that every
possible fairness intervention must fail.

% Sources: results/study2-compact-2026-10-21-aggregate.json; results/study2_rq1_aggregate_v1.json.
The linear follow-up addresses that weakness without changing the frozen representations
or participant roster. Its natural task-only reference reproduces the working ridge
predictor, with pooled SRCC 0.418 and normalised MAE 0.089. The separate outcome-only
selection analysis also selects a strongly regularised ridge in every outer fold.
Its mean fold SRCC of 0.567 and mean fold MAE of 0.081 describe that selection procedure;
they cannot be compared directly with a pooled OOF correlation as a paired improvement.
A stronger point estimate is not a substitute for a common estimand and uncertainty
analysis. This distinction follows the broader small-sample evaluation problem:
cross-validation estimates can vary substantially~\cite{Varoquaux2018}, and selection
on noisy development estimates can bias apparent performance~\cite{CawleyTalbot2010}.

Prediction validity remains conditional on the rating. The MIT outcome is a crowd-perceived
hireability judgement, not employment success. Ordinal Krippendorff alpha of 0.232 among
individual raters shows a substantial measurement limitation. The strong agreement of the
EM aggregate with an unweighted session mean indicates agreement between aggregation
methods, not that either is an objective criterion. Measurement choices determine what
the subsequent audit can establish~\cite{JacobsWallach2021}. Interview criterion validity
depends on the construct and later outcome measured~\cite{Wingate2025ConstructValidity};
neither corpus here contains a job-performance criterion.

\section{What the proxy penalty changes}

The original experiment supplies one clear component result. Adding W1 to the matched
visual or audio sampler lowers predictive utility. The incremental SRCC changes are
$-0.038$ [$-0.072$, $-0.005$] for visual proxies and $-0.047$
[$-0.091$, $-0.002$] for audio proxies. Original-scale MAE also worsens relative to
each sampling-only control. The matched design therefore prevents a combined
sampling-and-penalty change from being attributed to W1 alone. The random proxy and
adversarial reweighting comparisons further show why an arbitrary regulariser cannot
automatically be credited with meaningful group structure.

% Source: results/study2-compact-2026-10-21-aggregate.json.
The working linear model makes the mechanism question more informative. At strength
0.02, visual natural-exposure W1 reduces mean training proxy distance from 0.036 to
0.009 and mean held-out distance from 0.045 to 0.035. Audio natural exposure also reduces
held-out proxy W1, from 0.087 to 0.068. Thus the intervention is active in the observed
training and test distributions. The empirical interpretation no longer rests on
assuming that a numerically small coefficient must have been ineffective.

This response is accompanied by a trade-off. The visual natural-exposure SRCC change is
$-0.030$ [$-0.069$, 0.008], and the audio natural-exposure change is $-0.044$
[$-0.080$, $-0.008$]. Error changes are smaller and less resolved. A reduction in
proxy-score distance consequently does not establish a beneficial joint change in
ranking and error. The complete strength grid is more informative than a single favourable
point: the strongest audio penalty has considerably weaker ranking, while the smallest
penalty changes little. Thirteen highest-strength fits reach the declared iteration limit,
so these points are bounded numerical outcomes rather than evidence of converged optima.

The two studies also evaluate different balancing mechanisms. The original uses cyclic
group-balanced minibatches, whereas the linear follow-up uses the corresponding full-fold
group-weighted task objective. Its balanced and natural controls have nearly identical
rounded utility, but this does not erase the original sampling result or establish that
stochastic exposure is irrelevant. The shared lesson is that the penalty must be compared
with its own matched task-only objective, optimiser and exposure policy.

\section{Why proxy response does not establish audited fairness}

% Source: results/study2-compact-2026-10-21-aggregate.json.
The compact visual intervention reduces the post-session perceived-gender W1 point
estimate from 0.049 to 0.029. Its paired gain of approximately 0.020 has interval
[$-0.005$, 0.025]. This is a plausible observed transfer to that partition, but the
data do not resolve a reduction at the population level. The other perceived axes
give a different pattern: racialised-appearance point gains are negative, and skin-tone
gains are small, with all primary intervals spanning zero. A model can narrow one
learned partition while leaving another unchanged or increasing its score distance.

This is the missing implication identified in proxy-fairness research. Surrogate
constraints require assumptions about their relationship with the groups of
interest~\cite{Gupta2018ProxyFairness}. The present clustering and score penalty do not
supply those assumptions automatically. Nor does a lower W1 answer the error-balance
question. Score-distribution parity and prediction errors are distinct properties of a
regression model~\cite{Petersen2023RiskScores}; neither alone defines the fairness of a
future hiring decision.

The original perceived-attribute comparison is consistent with that caution. All six
primary neural-model disparity-gain intervals include zero, as do the matched incremental
penalty contrasts. Small distances also occur for weak mean predictors. A nearly constant
score can make two groups look similar while failing to predict the target. The linear
follow-up is valuable because it demonstrates proxy response on a stronger predictor,
yet it still does not establish a general improvement across the audited attributes.

The observer source bounds the interpretation. Perceived gender presentation, racialised
appearance and skin tone are responses to a particular recording, not measurements of
self-identified attributes. The 57 pre/post annotation pairs show stable gender codes but
only 31 exact appearance matches and 26 exact skin-tone matches. These are repeated
conditions rated by one source, not independent raters. Evidence on contextual subjectivity
in skin-tone annotation reinforces the need for caution~\cite{Barrett2023}; it cannot
supply a missing reliability estimate for this observer. Skin tone remains a five-point
ordinal field, and the maximum pairwise score distance does not test an ordinal trend.

The complete post audit is primary because it includes all 69 participants, not because
a post rating is intrinsically more valid. The 57-person pre sensitivity has a different
support pattern and twelve missing annotations whose missingness is not assumed random.
A positive secondary point or a lower interval endpoint rounding to zero cannot replace
the uncertain primary result. Bootstrap precision does not create more independent
people or recover unobserved rater disagreement.

\section{Representation and proxy meaning}

Visual and audio features occupy different measurement boundaries. MIT visual inputs are
candidate-face crops, while its audio contains both interviewer and participant.
The audio clusters may reflect room, channel, timing and interviewer behaviour as well as
candidate speech. Calling them candidate-voice or speaker-identity groups would overstate
what the representation measures.

The visual partition agrees exactly with this observer's post-session
gender-presentation codes in each held-out fold. Training-participant-only alignment
also permits a pooled description because all visual partitions are identical up to
label permutation. It yields the same perfect association across the 69 OOF
participants. This is a partition-association result within one corpus, not a general
demographic classifier or evidence of gender identity. It nevertheless makes an
important distinction visible: training without supplied gender labels does not mean
that the visual features contain no gender-presentation information. Audio's low
fold-local gender association and changing partitions show that the same interpretation
cannot simply be assigned to every modality.

Neutral k-means identifiers also have no cross-fold identity. An unaligned concatenation
can manufacture an apparent disagreement even when partitions are identical up to a
permutation. The held-out manipulation analysis therefore keeps fold-specific groups.
It does not use the invalid stitched grouping to justify demographic meaning or to
calculate a pooled proxy effect. Score disparities calculated directly from perceived
annotations remain separate from this cluster-number issue.

The original inference-only masks provide a further bounded explanation. Removing the
training-standardised audio summary lowers ERM SRCC in every seed, by 0.325 on average,
whereas removing video or text sometimes improves error. These are dependence diagnostics
of one weak fitted model, not causal evidence about human communication or the value of
retraining with fewer modalities. Temporal averaging discards order, and zero-mask
inference may create inputs unlike those seen in training. Such diagnostics motivate
simple competitive references rather than a claim that visual behaviour or language is
generally uninformative.

\section{What the independent corpus adds}

% Source: results/avi-baseline-2026-11-02-reconciled.json.
AVI-Personality adds independent participants and a different rating source
~\cite{Zhang2026AVI}. Its official test contains 130 people, and the selected ridge
predictor has original-scale MAE 0.419 against 0.479 for the training-target mean.
The approximately 12.6\% lower error shows useful predictive signal in its audio/video
features on this test. It does not imply that its raw MAE is better than MIT's:
the target scale, recruiter judgement, interview format and visual input boundary differ.

The group-error audit adds another source distinction. Self-reported male respondents
have test MAE 0.452 and female respondents 0.385, a descriptive signed gap of 0.067
rating points. One prefer-not-to-say response and the absence of a non-binary or
third-gender test respondent limit the categories that can be estimated. Those source
categories cannot be collapsed or moved between splits to create support. The positive
lower endpoint of a bootstrap interval for an absolute gap does not establish a nonzero
signed population difference; the reported point gap is therefore kept descriptive.

This external ridge analysis does not contain a W1 treatment/control comparison.
Consequently, it supports independent interview-score prediction and group-error
description rather than confirmation that the proxy penalty transports. That boundary
is central to RQ4. An external dataset does not, by itself, convert an exploratory MIT
finding into a confirmed general method effect. A method comparison formulated after
the external baseline and audit have been inspected is also exploratory, even though
its participants are independent of MIT.

The visual representation further limits an exact replication claim. AVI features
average sampled full frames, whereas MIT uses operationally selected face crops.
The independent release supplies audio/video but no transcripts. These differences can
affect both prediction and the proxy partition. They are part of the scope of the
external procedure, not a reason to pool results as if the same outcome and inputs had
been observed twice.


\section{The matched comparison on independent participants}

% Source: results/avi-compact-reconciled-v1.json.
The matched fifteen-arm comparison supplies the treatment-versus-control evidence the
baseline could not. Its primary contrast, the visual natural-exposure penalty at the
strength that served as the main study's primary setting, changes original-scale error by
0.002 rating points, rank correlation by $-0.018$ and the audited gender score gap by
0.001. Every interval contains zero. The independent corpus therefore reproduces the
main study's central outcome rather than correcting it: at the primary strength, proxy
regularisation leaves both utility and audited disparity where it found them.

The more informative result is the separation between training and held-out proxy response.
The penalty succeeds completely on its own optimisation target, collapsing the training
visual proxy gap from 0.021 to below 0.001 at the strongest setting. Almost none of that
reduction survives on held-out participants, where the same arm leaves the gap at 0.031
against a control value of 0.032. The failure is therefore one of generalisation, not of
optimisation: the constraint is satisfied on the participants used to fit it and does not
describe the participants it is then applied to. This distinction matters for interpreting
the main study, where the same pattern appears in weaker form, and it is not visible from
utility and disparity summaries alone.

The audio proxy behaves differently and is the more instructive case. Its constraint does
generalise, reducing the held-out proxy gap from 0.074 to 0.023. If reducing a proxy-group
score gap were sufficient to reduce an audited group gap, this is the arm in which the
effect should appear. It does not: the corresponding disparity gain is 0.012 with interval
$[-0.006, 0.029]$, while original-scale error rises by 0.034 and rank correlation falls by
0.178, both with intervals excluding zero. A proxy constraint that transfers successfully
to new participants thus purchased a substantial utility loss without a demonstrable
fairness return. This is direct evidence that the unsupervised partition is not a usable
stand-in for the attribute the audit measures, and it is stronger evidence than a null
obtained from a constraint that never transferred at all.

Three of the thirty-six intervals across the twelve penalised arms exclude zero on the
disparity measure. None is the primary contrast, no correction was applied for examining
twelve arms on three measures, and two of the three have lower endpoints within 0.0001 of
zero. The remaining one occurs at the strongest penalty, where ranking utility is also
declining. Treating any of these as a mitigation result would be selecting an outcome
after seeing it. They are reported because the complete grid is reported, and they are
interpreted as consistent with multiplicity rather than with a located effect.

One structural feature bounds all of these comparisons. Validation-selected shrinkage leaves
predicted scores with a standard deviation of 0.081 against an outcome standard deviation
of 0.149, before any penalty is applied. A score-disparity penalty acts on the separation
between two group prediction distributions, so a predictor compressed to roughly half the
dispersion of its target offers correspondingly little separation to remove. The regularisation
strength was selected at the upper limit of its search range in every fold of every study
reported here, which indicates a shrinkage-dominated regime rather than a boundary chosen
among competitive alternatives. The absence of disparity movement is consequently
consistent with two explanations that this evidence cannot separate: that proxy
regularisation does not reduce audited disparity, or that the scorers available to it
were too compressed for any such reduction to be expressed. The second explanation is
reinforced by the main study, where the learned configurations did not reach the utility of
the deterministic ridge floors, and by the random-proxy control, which performed as well as
the structured proxies.

\section{Answering the research questions}

\textbf{RQ1, whether a participant-safe benchmark supports a competitive scorer.} Partly,
and less well than intended. Correct pairing, a reliable-enough aggregated target and a
participant-grouped evaluation were established, and the deterministic floors demonstrate
recoverable signal, reaching rank correlation 0.424 and original-scale error 0.535 on MIT
and 0.473 and 0.419 on the independent corpus. The learned configurations did not match
those floors on any of the four utility measures. The benchmark is therefore sound as a
measurement instrument while the learned scorer built on it is not competitive, and that
gap conditions every answer that follows.

\textbf{RQ2, whether proxy regularisation measurably changes held-out proxy-score
distributions beyond balanced sampling.} Only partly, and the answer depends on the proxy.
In training the penalty controls the proxy gap completely, on both corpora. On held-out
participants the visual constraint barely transfers, while the audio constraint transfers
substantially. Balanced sampling alone did not reproduce the training effect. The mechanism
is therefore real but unreliable, and its reliability is a property of the specific
representation partitioned rather than of the method.

\textbf{RQ3, whether proxy change is accompanied by changes in perceived-group disparity at
an explicit utility cost.} No. On MIT the paired disparity intervals include zero, and on
the independent corpus the primary contrast does likewise. Where a proxy constraint did
transfer, the utility cost materialised without the disparity benefit. The utility cost is
thus demonstrable and the fairness return is not, which is the opposite of the trade-off the
method is intended to offer.

\textbf{RQ4, whether the evidence transports to an independent corpus.} The null transports;
a mitigation effect does not, because none was established to transport. The independent
comparison reproduces the main study's outcome under a different outcome construct, rating
source and visual boundary, which strengthens the negative finding. It does not convert any
exploratory result into a confirmed general effect, and the comparison was itself specified
after the external baseline had been inspected.

\section{Internal and statistical validity}

Participant isolation is essential in both corpora. The corrected MIT mapping keeps
each person's pre and post recording in one fold; AVI keeps all six responses in one
official participant split. Cross-fitting nevertheless reuses participants across training
sets, so folds and training seeds are not independent experimental samples. Dependence
also complicates general cross-validation uncertainty~\cite{BengioGrandvalet2004}.

The participant bootstrap conditions on the fitted models and observed splits. Its
paired draws preserve configuration matching, but it does not estimate the variability
of retraining and selecting a model on a fresh sample. Intervals are pointwise rather
than simultaneous guarantees across the full strength, proxy and audit grid. Fairness
measurement itself can have finite-sample bias~\cite{Lum2022}; neither a nonnegative
distance nor an absolute-gap interval is a general null test.

The strongest-penalty numerical terminations, mixed-speaker MIT audio, full-frame AVI
inputs, source-label missingness and low rating agreement are substantive limits.
The input and label separation protects the declared fit from direct audit selection;
it cannot remove bias already encoded upstream or validate the target against
employment outcomes. The reciprocal source correction for two MIT post sessions is
likewise a prerequisite for valid participant analysis rather than an independent
contribution to mitigation.

\section{Practical limits and further evidence}

The research audits scores in mock-interview datasets. It does not construct or
validate a hiring tool. Outcome prediction, score-distribution similarity and observed
group error must remain separate from job-performance validity, discrimination
experienced by applicants and compliance conclusions.

The most informative further evidence would address the limits exposed here: independently
documented perceived-attribute judgements, better outcome reliability, and an external
matched intervention comparison with a credible predictor. None can be replaced by
adding more seeds, bootstrap draws or weak model variants. Candidate-speech separation
and temporal modelling are also distinct experiments rather than retrospective repairs
of these results. The supported conclusion is an active proxy mechanism with a
conditional predictive cost and uncertain audited transfer, alongside an independent
demonstration of interview-score prediction.

\end{spacing}
